% ============================================================================
% LANGENTWURF LE-021: Hilfe holen
% Profil: S (Senioren 65+) | Tier: 2 (Standard)
% ============================================================================

\documentclass[11pt,a4paper]{article}
\usepackage[utf8]{inputenc}
\usepackage[T1]{fontenc}
\usepackage[ngerman]{babel}
\usepackage{geometry}
\usepackage{fancyhdr}
\usepackage{lastpage}
\usepackage{xcolor}
\usepackage{tcolorbox}
\usepackage{enumitem}
\usepackage{longtable}
\usepackage{hyperref}
\usepackage{amssymb}

\geometry{left=2.5cm, right=2.5cm, top=2.5cm, bottom=2.5cm}
\definecolor{fach}{HTML}{b35c00}
\definecolor{gray}{HTML}{666666}
\definecolor{successgreen}{HTML}{2a7a4b}

\tcbuselibrary{skins,breakable}
\newtcolorbox{scorebox}{ colback=white, colframe=fach, fonttitle=\bfseries, title=Didaktik-Score, breakable }
\newtcolorbox{infobox}[1][]{ colback=fach!5, colframe=fach, fonttitle=\bfseries, title=#1, breakable }

\pagestyle{fancy}
\fancyhf{}
\fancyhead[L]{\small\textcolor{gray}{Digitale Grundlagen -- 021 Hilfe holen}}
\fancyhead[R]{\small\textcolor{gray}{LE Tier 2 / Profil S}}
\fancyfoot[C]{\small\thepage{} / \pageref{LastPage}}
\renewcommand{\headrulewidth}{0.4pt}

\begin{document}

\begin{titlepage}
    \centering
    \vspace*{2cm}
    {\Large\textcolor{gray}{Digitale Grundlagen -- Senioren 65+}}\\[2cm]
    {\Huge\textcolor{fach}{\textbf{Langentwurf}}}\\[0.5cm]
    {\large Tier 2 | Profil S}\\[2cm]
    {\LARGE\textbf{021 -- Hilfe holen: Wer kann mir helfen?}}\\[0.5cm]
    {\large Modul E: Hilfe \& Selbstständigkeit}\\[2cm]
    \begin{tabular}{rl}
        \textbf{Zeitrahmen:} & 45--60 Minuten \\
        \textbf{Vorkenntnisse:} & Probleme lösen (020) \\
    \end{tabular}
    \vfill
    {\normalsize Stand: \today}
\end{titlepage}

\tableofcontents
\newpage

\section{Bedingungsanalyse}

\subsection{Ausgangssituation}
\begin{itemize}
    \item Scheu, um Hilfe zu bitten
    \item Nicht wissen, wer der richtige Ansprechpartner ist
    \item Angst vor hohen Kosten (``Reparatur'')
    \item Familie wird oft überfordert
    \item Apple Support unbekannt
\end{itemize}

\subsection{Typische Probleme}
\begin{itemize}
    \item ``An wen wende ich mich bei Problemen?''
    \item ``Kostet das etwas?''
    \item ``Wie erkläre ich mein Problem?''
    \item ``Ist der Apple Store für mich zuständig?''
\end{itemize}

\section{Sachanalyse}

\subsection{Kernkonzepte}
\begin{enumerate}
    \item \textbf{Vertrauensperson:} Familie, Freunde -- erste Anlaufstelle
    \item \textbf{Apple Support:} Offizieller Support (Telefon, Chat, Store)
    \item \textbf{VHS/Seniorenkurse:} Kurse für Anfänger in der Nähe
    \item \textbf{Problem beschreiben:} Was genau geht nicht? Seit wann?
\end{enumerate}

\subsection{Alltagsanalogien}
\begin{tabular}{|l|p{8cm}|}
\hline
\textbf{Begriff} & \textbf{Analogie} \\
\hline
Vertrauensperson & Nachbar, der bei kleinen Reparaturen hilft \\
Apple Support & Kundenservice wie bei der Telekom \\
Apple Store & Fachgeschäft mit Beratung \\
VHS-Kurs & Volkshochschule, aber für Computer \\
\hline
\end{tabular}

\section{Didaktische Analyse}

\subsection{Stellung in der Reihe}
Ergänzt Einheit 020. Nach der Selbsthilfe: Wo bekomme ich Hilfe?

\subsection{Didaktische Reduktion}
\textbf{Ausgeklammert:} Kostenpflichtige Reparaturen, Drittanbieter-Support.

\textbf{Fokus:} Drei Anlaufstellen kennen, Problem richtig beschreiben.

\section{Lernziele}

\begin{tabular}{|c|p{7cm}|c|c|}
\hline
\textbf{Nr.} & \textbf{Die Teilnehmer können...} & \textbf{Stufe} & \textbf{Niveau} \\
\hline
FZ1 & drei Anlaufstellen für Hilfe nennen & Wissen & Basis \\
FZ2 & ihr Problem klar beschreiben (Was? Seit wann?) & Anwenden & Basis \\
FZ3 & den Apple Support kontaktieren & Anwenden & Basis \\
FZ4 & eine Vertrauensperson festlegen & Anwenden & Basis \\
FZ5 & lokale Angebote (VHS) kennen & Wissen & Vertiefung \\
\hline
\end{tabular}

\section{Methodische Analyse}

\begin{infobox}[Die drei Anlaufstellen]
\begin{enumerate}
    \item \textbf{Vertrauensperson:} Familie, Freund, Nachbar -- bei kleinen Fragen
    \item \textbf{Apple Support:} Bei technischen Problemen -- Telefon: 0800 6645 451 (kostenlos)
    \item \textbf{VHS/Seniorenkurse:} Zum Lernen und Üben -- Kurse in deiner Stadt
\end{enumerate}
\end{infobox}

\section{Verlaufsplanung}

\begin{longtable}{|p{1cm}|p{2cm}|p{5cm}|p{3cm}|p{2cm}|}
\hline
\textbf{Zeit} & \textbf{Phase} & \textbf{Inhalt} & \textbf{TN-Aktivität} & \textbf{Material} \\
\hline
0--10 & Einstieg & ``Wen rufen Sie, wenn die Heizung nicht geht?'' Parallele & Berichten & -- \\
\hline
10--20 & Vertrauensperson & Wer in der Familie/Freundeskreis kann helfen? & Nachdenken, aufschreiben & S.1 \\
\hline
20--35 & Apple Support & Telefonnummer, Website, Chat, Store vorstellen & Nummer notieren & S.1 \\
\hline
35--45 & Problem beschreiben & Übung: ``Mein Problem ist...'' formulieren & Üben & S.2 \\
\hline
45--55 & Lokale Angebote & VHS, Seniorenkurse in der Region & Infos erhalten & S.2-3 \\
\hline
\end{longtable}

\section{Lösungen}

\textbf{Apple Support kontaktieren:}
\begin{itemize}
    \item \textbf{Telefon:} 0800 6645 451 (kostenlos aus Deutschland)
    \item \textbf{Website:} support.apple.com/de-de
    \item \textbf{App:} ``Apple Support'' im App Store
    \item \textbf{Store:} Termin im Apple Store vereinbaren
\end{itemize}

\textbf{Problem richtig beschreiben:}
\begin{itemize}
    \item \textbf{Was} geht nicht? (``Die Mail-App öffnet nicht'')
    \item \textbf{Seit wann?} (``Seit gestern'')
    \item \textbf{Was habe ich versucht?} (``Neustart hat nicht geholfen'')
    \item \textbf{Fehlermeldung?} (Falls vorhanden, notieren oder fotografieren)
\end{itemize}

\textbf{Lokale Angebote finden:}
\begin{itemize}
    \item Volkshochschule (VHS): www.vhs.de
    \item Seniorenbüro der Stadt
    \item Mehrgenerationenhäuser
    \item Bibliotheken (oft kostenlose Kurse)
\end{itemize}

\section{Didaktik-Score}

\begin{scorebox}
\begin{tabular}{|c|l|c|c|p{3.5cm}|}
\hline
\# & Dimension & Gew. & Score & Notiz \\
\hline
1 & Differenzierung & 15\% & 9/10 & Drei Stufen: Familie, Apple, VHS \\
2 & Taxonomie & 10\% & 9/10 & Problem beschreiben = Anwenden \\
3 & Alltagsrelevanz & 10\% & 10/10 & Jeder braucht Hilfe \\
4 & Aktivierung & 10\% & 9/10 & Vertrauensperson festlegen \\
5 & Scaffolding & 10\% & 9/10 & Klare Struktur \\
6 & Feedback & 10\% & 9/10 & Rollenspiel mit Ermutigung \\
7 & Sprachliche Klarheit & 10\% & 9/10 & Einfach \\
8 & Motivation & 10\% & 9/10 & Sicherheit geben \\
9 & Barrierefreiheit & 10\% & 9/10 & Große Schrift \\
10 & Praktikabilität & 5\% & 9/10 & 55 Min gut \\
\hline
\multicolumn{3}{|r|}{\textbf{GESAMT:}} & \textbf{9.10/10} & \textcolor{successgreen}{\textbf{FREIGABE}} \\
\hline
\end{tabular}
\end{scorebox}

\end{document}
